% =============================================================================
% Chapter 3: Methodology (8-10 pages)
% =============================================================================

\section{Methodology}
\label{sec:methodology}

This chapter presents the technical core of the thesis: the QUBO
formulations of modularity maximization, the GNN architecture, the
loss function with regularization, and the training procedure.

% -----------------------------------------------------------------------------
\subsection{QUBO formulation of modularity}
\label{sec:method_qubo}

\subsubsection{Newman's modularity}
\label{sec:method_modularity}

The modularity of a partition $C: V \to \{1, \ldots, k\}$ of an
undirected graph $G = (V, E)$ is defined~\cite{newman2006modularity}
as
\begin{equation}
    M(C) = \frac{1}{2m} \sum_{v, w \in V} B_{vw} \, \delta(C_v, C_w),
    \label{eq:modularity}
\end{equation}
where $B_{vw} = A_{vw} - k_v k_w / (2m)$ is the entry of the
\emph{modularity matrix}, $A_{vw}$ is the adjacency matrix entry
between vertices $v$ and $w$, $k_v = \sum_w A_{vw}$ is the degree of
vertex $v$, $m = \frac{1}{2}\sum_{v, w} A_{vw}$ is the total number
of edges, and $\delta$ is the Kronecker delta.

The intuition behind~\eqref{eq:modularity} is straightforward:
$B_{vw}$ measures the difference between the actual number of edges
between $v$ and $w$ and the expected number under a random graph
with the same degree sequence. Modularity $M(C) > 0$ thus indicates
that the partition groups together vertices that are more densely
connected than expected by chance, while $M(C) < 0$ suggests an
``anti-community'' structure. Maximizing $M$ has become the standard
objective for community detection, both in classical algorithms such
as Louvain~\cite{blondel2008louvain} and in optimization-based
formulations~\cite{wang2024quantum}.

A useful property of the modularity matrix that follows directly
from the definition is
\begin{equation}
    \sum_{w \in V} B_{vw} = 0 \quad \text{for all } v \in V,
    \label{eq:B_row_sum}
\end{equation}
which we will use in the QUBO derivation
(Sections~\ref{sec:method_qubo_k2}
and~\ref{sec:method_qubo_multi}).

\subsubsection{QUBO general form}
\label{sec:method_qubo_general}

A Quadratic Unconstrained Binary Optimization problem~\cite{lucas2014ising}
takes the form
\begin{equation}
    \min_{\xb \in \{0,1\}^N} \, \xb^T \Q \xb,
    \label{eq:qubo_general}
\end{equation}
where $\Q \in \mathbb{R}^{N \times N}$ is a symmetric matrix
encoding both the linear and quadratic interactions among the
variables. The structure of $\Q$ has a direct combinatorial
interpretation. Since $x_i^2 = x_i$ for binary variables, the
diagonal entry $Q_{ii}$ acts as a linear coefficient on $x_i$,
while the off-diagonal entry $Q_{ij}$ contributes $Q_{ij} x_i x_j$
to the objective whenever both $x_i$ and $x_j$ are set to one.
Without loss of generality $\Q$ may be assumed upper-triangular,
since pairs $(Q_{ij}, Q_{ji})$ can always be combined into a single
entry $Q_{ij} + Q_{ji}$.

QUBO is NP-hard in general~\cite{lucas2014ising}, and most
combinatorial optimization problems of practical interest~--~maximum
cut, graph coloring, set cover, and many others~--~admit a QUBO
formulation. The same is true for community detection by modularity
maximization, as developed in the next subsections.


\subsubsection{Modularity as QUBO for $k = 2$}
\label{sec:method_qubo_k2}

For two communities, a single binary variable per vertex suffices.
Let $x_v = 1$ if vertex $v$ belongs to community A and $x_v = 0$
if it belongs to community B. The Kronecker delta in the modularity
expression~\eqref{eq:modularity} can then be rewritten as
\begin{equation}
    \delta(C_v, C_w) = x_v x_w + (1 - x_v)(1 - x_w),
    \label{eq:delta_k2}
\end{equation}
since the right-hand side is $1$ when $x_v = x_w$ and $0$ otherwise.

Substituting~\eqref{eq:delta_k2} into~\eqref{eq:modularity} gives,
after expanding the product,
\begin{equation*}
    M(C) = \frac{1}{2m} \sum_{v,w} B_{vw} \bigl( 2 x_v x_w - x_v - x_w
    + 1 \bigr).
\end{equation*}
The terms involving a single $x$ vanish by the row-sum
property~\eqref{eq:B_row_sum}, since $\sum_w B_{vw} = 0$ implies
$\sum_{v,w} B_{vw} \, x_v = 0$, and the constant term $\sum_{v,w}
B_{vw}$ also vanishes. The result is
\begin{equation}
    M(C) = \frac{1}{m} \sum_{v,w} B_{vw} \, x_v x_w.
    \label{eq:modularity_k2_quadratic}
\end{equation}
Maximizing $M$ is therefore equivalent to minimizing
\begin{equation}
    H_2(\xb) = -\sum_{v,w} B_{vw} \, x_v x_w = \xb^T (-\B) \, \xb,
    \label{eq:qubo_k2}
\end{equation}
which is a QUBO of the form~\eqref{eq:qubo_general} with
$\Q = -\B$ (up to a positive scaling factor that does not affect
the optimizer). The problem size is $N = n$, and the formulation
contains only quadratic terms; no slack variables or constraint
penalties are needed.


\subsubsection{Modularity as QUBO for $k > 2$}
\label{sec:method_qubo_multi}

When more than two communities are required, a single binary
variable per vertex no longer suffices. The standard
remedy~\cite{negre2020detecting, wang2024quantum} is to introduce
$k$ binary variables per vertex and constrain exactly one of them
to be active.

\paragraph{One-hot encoding.}
For each vertex $v$ and community $c \in \{1, \ldots, k\}$, define
\begin{equation*}
    x_{v,c} = \begin{cases} 1 & \text{if vertex } v \text{ is
    assigned to community } c, \\ 0 & \text{otherwise.} \end{cases}
\end{equation*}
A valid community assignment satisfies the one-hot constraint
\begin{equation}
    \sum_{c=1}^{k} x_{v,c} = 1 \qquad \text{for all } v \in V.
    \label{eq:one_hot}
\end{equation}
The modularity~\eqref{eq:modularity} expressed in these variables
becomes
\begin{equation}
    M(C) = \frac{1}{2m} \sum_{c=1}^{k} \sum_{v,w \in V} B_{vw} \,
    x_{v,c} x_{w,c},
    \label{eq:modularity_onehot}
\end{equation}
where the inner sum gathers all pairs of vertices assigned to
community $c$.

\paragraph{Constraint penalty.}
The QUBO form is by definition unconstrained, so the one-hot
constraint~\eqref{eq:one_hot} cannot be enforced directly. Instead
it is incorporated as a soft penalty term~\cite{lucas2014ising}:
\begin{equation*}
    P \sum_{v \in V} \left( \sum_{c=1}^{k} x_{v,c} - 1 \right)^2,
\end{equation*}
where $P > 0$ is a penalty coefficient chosen large enough to make
constraint violations more costly than any modularity gain. The
penalty is zero when the one-hot constraint is satisfied and
positive otherwise.

\paragraph{Full formulation.}
Combining the negated modularity with the constraint penalty, the
QUBO objective for $k$ communities, following formula~5 of Wang et
al.~\cite{wang2024quantum}, is
\begin{equation}
    H_k(\xb) = -\frac{1}{2m} \sum_{c=1}^{k} \sum_{v,w \in V} B_{vw} \,
    x_{v,c} x_{w,c} \;+\; P \sum_{v \in V} \left(
    \sum_{c=1}^{k} x_{v,c} - 1 \right)^2.
    \label{eq:qubo_multi}
\end{equation}

\paragraph{Linearization.}
To put~\eqref{eq:qubo_multi} into the standard QUBO
form~\eqref{eq:qubo_general}, the double index $(v, c)$ is
linearized into a single index
\begin{equation}
    i = (v - 1) \cdot k + c,
    \label{eq:linearization}
\end{equation}
giving a flat vector $\mathbf{y} \in \{0, 1\}^N$ with $N = nk$. The
quadratic form $\mathbf{y}^T Q \mathbf{y}$ that
reproduces~\eqref{eq:qubo_multi} has the following block
structure, given in the upper-triangular convention
(Section~\ref{sec:method_qubo_general}). Diagonal entries
$Q_{ii}$ combine the modularity contribution $-B_{vv} / (2m)$
with the linear part of the constraint penalty $-P$. Off-diagonal
entries $Q_{ij}$ with $i < j$ fall into one of two categories: a
\emph{modularity entry} $Q_{ij} = -B_{vw} / m$ when indices
correspond to different vertices but the same community, and a
\emph{constraint entry} $Q_{ij} = 2P$ when indices correspond to
the same vertex but different communities. All remaining entries
are zero. The resulting matrix has size $nk \times nk$. (An
equivalent symmetric storage halves these off-diagonal values and
sets $Q_{ji} = Q_{ij}$; the implementation in this thesis uses
the upper-triangular convention throughout.)

\paragraph{Choice of $P$.}
The penalty $P$ must be large enough to enforce the one-hot
constraint but not so large that the modularity term is dominated
and the optimization landscape is flattened. In the experiments
that follow, $P = 1.5 \cdot \max_{v,w} |B_{vw}|$ is used, which
puts the constraint and modularity terms on comparable scales.


\subsubsection{QUBO and Ising equivalence}
\label{sec:method_qubo_ising}

The QUBO and Ising forms of an optimization problem are equivalent
under the affine change of variables $s_i = 2 x_i - 1$, which maps
$x_i \in \{0, 1\}$ to $s_i \in \{-1, +1\}$. Substituting into a
QUBO objective $\xb^T \Q \xb$ produces an Ising
Hamiltonian~\cite{lucas2014ising} of the form
\begin{equation*}
    H_{\mathrm{Ising}}(\mathbf{s}) = \sum_i h_i s_i + \sum_{i < j} J_{ij}
    s_i s_j + \mathrm{const},
\end{equation*}
where the local fields $h_i$ and couplings $J_{ij}$ are determined
by the original $\Q$. The two formulations have the same set of
minimizers up to the affine relabeling, so any solver designed for
one can be used for the other. This thesis works directly with the
QUBO form, but the equivalence is invoked when comparing with
hardware Ising solvers such as those of Wang et
al.~\cite{wang2024quantum}.

% -----------------------------------------------------------------------------
\subsection{GNN architecture}
\label{sec:method_gnn}

The neural component of the proposed method is a graph neural
network whose output, after softmax normalization, parameterizes
the relaxed QUBO assignment. The architecture follows the QIGNN
framework~\cite{qignn2026}, with the multi-class adaptations
described in Section~\ref{sec:method_ressage}.

\subsubsection{Message passing}
\label{sec:method_messagepassing}

Graph neural networks operate on graph-structured data by
maintaining a vector representation $h_v^{(\ell)} \in
\mathbb{R}^{d_\ell}$ for each vertex $v$ at each layer $\ell$, and
updating these representations layer by layer through messages
exchanged with neighboring vertices. The general
update~\cite{hamilton2017inductive, kipf2017semi} takes the form
\begin{equation}
    h_v^{(\ell+1)} = \mathrm{UPDATE}\bigl( h_v^{(\ell)}, \,
    \mathrm{AGG}(\{h_u^{(\ell)} : u \in \Ncal(v)\}) \bigr),
    \label{eq:message_passing}
\end{equation}
where $\Ncal(v)$ is the set of neighbors of $v$, $\mathrm{AGG}$ is
a permutation-invariant aggregation function over the neighbors,
and $\mathrm{UPDATE}$ is a learned transformation that combines the
aggregated neighbor information with the vertex's own current
representation. Different choices of $\mathrm{AGG}$ and
$\mathrm{UPDATE}$ give rise to different GNN architectures.


\subsubsection{GraphSAGE}
\label{sec:method_sage}

The architecture used in this thesis is built from
GraphSAGE~\cite{hamilton2017inductive} convolutional layers. A
GraphSAGE layer with aggregator $\mathrm{AGG}$ updates each vertex
representation through
\begin{equation}
    h_v^{(\ell+1)} = \sigma\!\left( W^{(\ell)} \left[
    h_v^{(\ell)} \,\Vert\, \mathrm{AGG}\bigl( \{ h_u^{(\ell)} : u
    \in \Ncal(v) \} \bigr) \right] \right),
    \label{eq:sage_update}
\end{equation}
where $\Vert$ denotes vector concatenation, $W^{(\ell)}$ is a
learned weight matrix, and $\sigma$ is a non-linear activation. The
two main aggregator choices for $\mathrm{AGG}$ are the mean
aggregator, which averages the neighbor representations
componentwise, and the max-pool aggregator, which applies a small
multilayer perceptron to each neighbor representation and then
takes the componentwise maximum.


\subsubsection{SAGEResBlock with dual aggregation}
\label{sec:method_resblock}

The basic building block of the QIGNN architecture is the
\texttt{SAGEResBlock}, which sums two parallel GraphSAGE branches
with different aggregators (the name follows the original QIGNN
code; the block performs branch-sum aggregation rather than a
classical identity-skip residual connection).
Given an input representation $h$, the block computes
\begin{equation}
    \mathrm{SAGEResBlock}(h) = \mathrm{LeakyReLU}\bigl(
    \mathrm{BN}_1(\mathrm{SAGE}_{\mathrm{mean}}(h)) +
    \mathrm{BN}_2(\mathrm{SAGE}_{\mathrm{pool}}(h)) \bigr),
    \label{eq:resblock}
\end{equation}
where $\mathrm{SAGE}_{\mathrm{mean}}$ and
$\mathrm{SAGE}_{\mathrm{pool}}$ are GraphSAGE layers with mean and
max-pool aggregators respectively, $\mathrm{BN}_1$ and
$\mathrm{BN}_2$ are batch normalization layers, and the addition
is performed elementwise. The motivation for the dual aggregation
is empirical: the QIGNN authors~\cite{qignn2026} report that
combining the two aggregators provides a richer neighborhood
summary than either aggregator alone, improving final solution
quality on combinatorial benchmarks.


\subsubsection{ResSAGE stack}
\label{sec:method_ressage}

The full network, called \texttt{ResSAGE} after the original
implementation, consists of one or more \texttt{SAGEResBlock}
layers followed by a final GraphSAGE layer that produces the output.
Throughout this thesis a single \texttt{SAGEResBlock} is used,
followed by dropout and a final layer that maps to the desired
output dimension. The output dimension and final activation depend
on the QUBO formulation:
\begin{itemize}
    \item For the bipartition formulation~\eqref{eq:qubo_k2}, the
    final layer outputs a single logit per vertex, followed by a
    sigmoid activation. The result is a vector $\pb \in [0,1]^n$
    representing the relaxed assignment to community A.

    \item For the multi-class formulation~\eqref{eq:qubo_multi},
    the final layer outputs $k$ logits per vertex, followed by a
    softmax along the community dimension. The result is a matrix
    $\Pb \in [0,1]^{n \times k}$ whose rows sum to one and represent
    a relaxed one-hot assignment over $k$ communities.
\end{itemize}
The replacement of the sigmoid head with a softmax over $k$ classes
is one of the architectural adaptations made in this work; the
original QIGNN architecture supports only single-bit-per-vertex outputs.


\subsubsection{Iterative refinement}
\label{sec:method_iterative}

A distinguishing feature of QIGNN relative to the simpler
PI-GNN~\cite{schuetz2022combinatorial} is the use of iterative
refinement during training. At training epoch $t$, the network is
given two inputs: a static initial vertex feature matrix $h^{(0)}$
(random embeddings concatenated with PageRank scores) and the
network's own output from the previous epoch, $h_0^{(t-1)}$. The
two are concatenated into a single input
\begin{equation}
    \tilde{h}^{(t)} = \bigl[ h^{(0)} \,\Vert\, h_0^{(t-1)} \bigr],
    \label{eq:iterative_input}
\end{equation}
which is then passed through the network. On the first epoch, the
previous output is initialized to zero, $h_0^{(0)} = \mathbf{0}$.
This feedback mechanism allows the network to revisit and improve
its earlier predictions rather than re-deriving them from scratch
at each epoch.


% -----------------------------------------------------------------------------
\subsection{Continuous relaxation and loss function}
\label{sec:method_loss}

\subsubsection{Continuous relaxation}
\label{sec:method_relaxation}

Both QUBO formulations developed in
Section~\ref{sec:method_qubo}~--~the bipartition $\xb^T (-\B) \xb$
and the multi-class form~\eqref{eq:qubo_multi}~--~minimize a
quadratic objective over a discrete domain $\{0, 1\}^N$. The
discrete domain admits no notion of gradient: there is no
infinitesimal direction in which a binary variable can be moved.
Direct application of gradient-based optimization is therefore
impossible.

The standard remedy~\cite{schuetz2022combinatorial, lucas2014ising}
is continuous relaxation. The discrete domain $\{0, 1\}^N$ is
replaced by the continuous hypercube $[0, 1]^N$, on which the
quadratic form $\pb^T \Q \pb$ is a polynomial of $\pb$ and is
therefore smooth and differentiable. The original discrete vertices
of the hypercube are still feasible points, but the optimizer is
allowed to traverse the interior of the cube during training. After
training, a feasible discrete solution is recovered by rounding,
discussed in Section~\ref{sec:method_postprocessing}.

The relaxation is exact for linear objectives, whose minimum on a
hypercube is always attained at a vertex, but in general lossy for
indefinite quadratic forms such as the QUBO objectives considered
here. The minimum of the relaxed problem may lie strictly inside
the hypercube, and rounding to a discrete vertex can introduce a
gap between the relaxed and discrete optima. The empirical
evaluation in Chapter~\ref{sec:experiments} measures the size of
this gap on a range of benchmarks.


\subsubsection{Differentiability through the GNN}
\label{sec:method_differentiability}

In the proposed method the relaxed assignment $\pb$ is not optimized
directly. Instead it is parameterized by the GNN: $\pb =
\pb(\theta)$, where $\theta$ are the GNN parameters. The training
loss is then a function of $\theta$:
\begin{equation*}
    \Lcal(\theta) = \pb(\theta)^T \, \Q \, \pb(\theta).
\end{equation*}
The loss is differentiable in $\theta$ because each component in
the chain $\theta \mapsto \pb \mapsto \Lcal$ is differentiable: the
GNN is a composition of matrix multiplications, additions, and
elementwise non-linearities (LeakyReLU, sigmoid, softmax), each of
which has a well-defined derivative; and the quadratic form is
polynomial in $\pb$. Backpropagation through this composite
function is performed automatically by the autograd
machinery~\cite{paszke2019pytorch} and yields the gradient
$\nabla_\theta \Lcal$ used by the optimizer.

This is the rationale for parameterizing the assignment through a
GNN rather than optimizing $\pb$ directly. The graph structure is
implicitly built into the GNN through message passing: vertices
that are close in the graph share information through their layers
and tend to receive similar embeddings, and hence similar relaxed
assignments. Direct gradient descent on $\pb \in \mathbb{R}^N$
exploits no such structure.


\subsubsection{QUBO loss}
\label{sec:method_qubo_loss}

For the bipartition formulation the loss is
\begin{equation}
    \Lcal_{\mathrm{bip}}(\theta) = \pb(\theta)^T \, (-\B) \,
    \pb(\theta),
    \label{eq:loss_bipartition}
\end{equation}
where $\pb(\theta) \in [0,1]^n$ is the GNN output after sigmoid
activation.

For the multi-class formulation the GNN produces a probability
matrix $\Pb(\theta) \in [0,1]^{n \times k}$ after softmax
normalization. To match the QUBO form~\eqref{eq:qubo_general}, the
matrix is flattened by the linearization~\eqref{eq:linearization}
into a vector $\pb_{\mathrm{flat}}(\theta) \in [0,1]^{nk}$, and the
loss takes the form
\begin{equation}
    \Lcal_{\mathrm{QUBO}}(\theta) = \pb_{\mathrm{flat}}(\theta)^T
    \, \Q \, \pb_{\mathrm{flat}}(\theta).
    \label{eq:loss_qubo}
\end{equation}


\subsubsection{Annealing penalty}
\label{sec:method_annealing}

Both losses~\eqref{eq:loss_bipartition} and~\eqref{eq:loss_qubo}
are augmented with a time-dependent penalty term that pushes the
relaxed solution toward the corners of the hypercube as training
progresses~\cite{schuetz2022combinatorial}. The penalty is
\begin{equation}
    \Lcal_{\mathrm{anneal}}(\theta, t) = \lambda(t) \sum_{i}
    p_i (1 - p_i),
    \label{eq:loss_anneal}
\end{equation}
where $p_i$ ranges over all entries of $\pb$ (or
$\pb_{\mathrm{flat}}$ in the multi-class case) and the schedule
$\lambda(t) = t / T_0$ with $T_0 = 10^4$ grows linearly with the
epoch number $t$. The product $p_i (1 - p_i)$ is maximized at
$p_i = 0.5$ and zero at $p_i \in \{0, 1\}$, so the penalty
effectively rewards confident, near-discrete assignments. The
linear schedule starts the training without bias toward any
particular point of the hypercube and gradually introduces
pressure toward the discrete vertices.


\subsubsection{Memory-efficient structured loss}
\label{sec:method_efficient}

Direct evaluation of the multi-class
loss~\eqref{eq:loss_qubo} as written would require materializing
the QUBO matrix $\Q \in \mathbb{R}^{nk \times nk}$ in memory. For
graphs of practical interest this is prohibitive: on the
Email-EU-core dataset ($n = 986$, $k = 42$), the matrix has
$N^2 \approx 1.72 \cdot 10^9$ entries, or approximately~$6.9$~GB
at single precision. A direct implementation is therefore unable
to handle even moderately large graphs.

The full $\Q$ matrix can be avoided by exploiting the block
structure described in Section~\ref{sec:method_qubo_multi}.
The QUBO objective~\eqref{eq:qubo_multi} decomposes into a
modularity term and a constraint penalty term. The modularity term
\begin{equation*}
    -\frac{1}{2m} \sum_{c=1}^{k} \sum_{v,w \in V} B_{vw} \,
    P_{v,c} P_{w,c}
\end{equation*}
sums over pairs $(v, w)$ of vertices that share a community $c$,
weighted by the modularity matrix entry $B_{vw}$. Pulling the
sum over $c$ inside, this is exactly the trace of $\Pb^T \B \Pb$:
\begin{equation}
    \Lcal_{\mathrm{mod}}(\Pb) = -\frac{1}{2m} \, \mathrm{tr}\bigl(
    \Pb^T \B \, \Pb \bigr).
    \label{eq:loss_mod_efficient}
\end{equation}
The constraint penalty
\begin{equation*}
    P \sum_{v \in V} \left( \sum_{c=1}^{k} P_{v,c} - 1 \right)^2
\end{equation*}
can be written compactly as the squared $\ell_2$ norm of the row
sums of $\Pb$ minus one:
\begin{equation}
    \Lcal_{\mathrm{constraint}}(\Pb) = P \, \bigl\| \Pb \mathbf{1}_k
    - \mathbf{1}_n \bigr\|_2^2,
    \label{eq:loss_constraint_efficient}
\end{equation}
where $\mathbf{1}_k$ and $\mathbf{1}_n$ are vectors of ones of the
appropriate dimensions. The total multi-class loss is then
\begin{equation}
    \Lcal_{\mathrm{QUBO}}(\theta) = \Lcal_{\mathrm{mod}}(\Pb(\theta))
    + \Lcal_{\mathrm{constraint}}(\Pb(\theta)),
    \label{eq:loss_total_efficient}
\end{equation}
which is mathematically equivalent to~\eqref{eq:loss_qubo}: the
block structure of $\Q$ described in
Section~\ref{sec:method_qubo_multi} factors exactly into the
modularity trace and the constraint norm term, so no constants or
cross-terms are lost.

The structured form requires only the modularity matrix
$\B \in \mathbb{R}^{n \times n}$ and the probability matrix
$\Pb \in \mathbb{R}^{n \times k}$, with no need to materialize
$\Q$. Memory complexity drops from $O(n^2 k^2)$ for the dense
QUBO representation to $O(n^2 + nk)$ for the structured
representation. A sparse implementation of $\Q$ would already
reduce the cost to $O(n^2 k + n k^2)$ by exploiting the block
structure; the structured form goes one factor further by
treating the modularity and constraint terms separately, and is
$k^2$ times smaller than the dense baseline when the modularity
term dominates. For the Email-EU-core example ($k=42$), this is
a reduction from approximately~$6.9$~GB (dense baseline) to
approximately~$4$~MB. Without this reformulation, training on
graphs with $n \sim 10^3$ and $k \sim 10^2$ would not be
feasible on commodity hardware.


\subsubsection{Output postprocessing}
\label{sec:method_postprocessing}

Training produces a relaxed assignment $\pb \in [0,1]^n$ for the
bipartition case or $\Pb \in [0,1]^{n \times k}$ for the multi-class
case. A discrete partition is recovered by rounding:
\begin{itemize}
    \item Bipartition: $C_v = 1$ if $p_v \geq 0.5$ and $C_v = 0$
    otherwise.
    \item Multi-class: $C_v = \arg\max_c P_{v,c}$.
\end{itemize}
The annealing penalty~\eqref{eq:loss_anneal} ensures that, by the
end of training, most entries of $\pb$ or $\Pb$ are close to either
$0$ or $1$, so the rounding step rarely changes the result
materially.

% -----------------------------------------------------------------------------
\subsection{Regularization for trivial collapse}
\label{sec:method_regularization}

\subsubsection{Trivial collapse problem}
\label{sec:method_collapse_problem}

The multi-class QUBO loss~\eqref{eq:loss_qubo} has a wide and
easily reachable local attractor corresponding to the trivial
assignment in which every vertex is placed in the same community
$c^*$, that is, $P_{v,c^*} = 1$ for all $v$ and $P_{v,c} = 0$ for
all $c \neq c^*$. In this state the modularity contribution
\begin{equation*}
    -\frac{1}{2m} \sum_{c=1}^{k} \sum_{v,w} B_{vw} P_{v,c} P_{w,c}
    = -\frac{1}{2m} \sum_{v,w} B_{vw}
\end{equation*}
vanishes by the row-sum property~\eqref{eq:B_row_sum}, and the
constraint penalty also vanishes since $\sum_c P_{v,c} = 1$ for
every vertex. The trivial assignment thus achieves a loss of
\emph{zero} --- which is \emph{worse} than any non-trivial valid
assignment that has $M(C) > 0$, since the loss at a valid
one-hot assignment equals $-M(C)$. The trivial state is therefore
not a global minimum, but it is a wide attractor in the loss
landscape: it requires no coordinated agreement between vertices
and is reached as soon as the softmax outputs are biased toward
any single community. Empirically the optimization converges to
this attractor in a substantial fraction of training runs
(Section~\ref{sec:exp_collapse}).


\subsubsection{Orthogonality regularization}
\label{sec:method_ortho}

To discourage trivial collapse, an orthogonality regularization
adapted from the supervised MinCutPool method~\cite{bianchi2020spectral}
(also used in DMoN~\cite{tsitsulin2020graph}) is added to the
multi-class loss:
\begin{equation}
    \Lcal_{\mathrm{ortho}}(\Pb) = \left\| \frac{\Pb^T \Pb}{\| \Pb^T
    \Pb \|_F} - \frac{I_k}{\sqrt{k}} \right\|_F^2,
    \label{eq:loss_ortho}
\end{equation}
where $\| \cdot \|_F$ is the Frobenius norm and $I_k$ is the
$k \times k$ identity matrix. The matrix $\Pb^T \Pb$ collects
inner products between the community columns of $\Pb$: its
diagonal entry $(c, c)$ is the squared mass of community $c$, and
its off-diagonal entry $(c, c')$ measures the overlap of
communities $c$ and $c'$. The minimum of
$\Lcal_{\mathrm{ortho}}$ is reached when $\Pb^T \Pb$ is proportional
to $I_k$, that is, when the community columns are mutually
orthogonal and of equal magnitude. This rules out trivial
collapse, which produces a $\Pb^T \Pb$ with mass concentrated in a
single diagonal entry.

A side effect of this regularization is that it also encourages
communities to be of approximately equal size, since equal column
norms appear in the optimal $\Pb^T \Pb$. This is a strong prior
that aligns with balanced benchmarks but conflicts with the
natural structure of imbalanced real-world graphs. The empirical
analysis in Section~\ref{sec:exp_regularization} characterizes the
resulting trade-off across graphs with different size
distributions.

\subsection{Hierarchical decomposition algorithm}
\label{sec:method_hierarchical}

The bipartition QUBO formulation~\eqref{eq:qubo_k2} converges
reliably and produces high-quality two-way splits, while the
multi-class formulation~\eqref{eq:qubo_multi} suffers from the
trivial collapse described in Section~\ref{sec:method_collapse_problem}
and is sensitive to the choice of regularization. A natural way to
exploit the strength of the bipartition formulation while still
addressing problems with $k > 2$ communities is to apply the
bipartition formulation \emph{recursively}: split the graph into
two parts, then split each part again, and continue until further
splitting no longer improves modularity.

This idea is similar in spirit to the multi-level scheme of
Louvain~\cite{blondel2008louvain}, which alternates between
greedy modularity gain at the vertex level and aggregation into
super-vertices, but the splitting step is performed by a neural
QUBO solver rather than greedy local search. To the best of the
author's knowledge, this combination has not been studied
previously.

\paragraph{Algorithm.}
The proposed procedure is summarized in
Algorithm~\ref{alg:hierarchical}. The input is a graph $G$, an
optional target number of communities $k_{\max}$, and a minimum
community size $n_{\min}$ below which a candidate community is
not split further. A queue holds the current set of candidate
subgraphs to be examined. Each candidate is split using the
bipartition QIGNN of Section~\ref{sec:method_qubo_k2}, and the
split is accepted if and only if it produces a positive modularity
gain on the original graph. Accepted splits add two children to
the queue; rejected splits or splits of subgraphs below $n_{\min}$
finalize the current subgraph as a community.

\begin{algorithm}[H]
\caption{Hierarchical QIGNN}
\label{alg:hierarchical}
\begin{algorithmic}[1]
\Require Graph $G = (V, E)$, target $k_{\max}$, minimum size $n_{\min}$
\State $\Cscr \gets \emptyset$ \Comment{final communities}
\State $Q \gets [V]$ \Comment{queue of candidate vertex sets}
\While{$Q$ is non-empty \textbf{and} $|\Cscr| + |Q| < k_{\max}$}
    \State $S \gets$ pop largest set from $Q$
    \If{$|S| \le n_{\min}$}
        \State $\Cscr \gets \Cscr \cup \{S\}$
        \State \textbf{continue}
    \EndIf
    \State $(S_1, S_2) \gets \textsc{BipartitionQIGNN}(G[S])$
    \State $\Delta M \gets M(\Cscr \cup \{S_1, S_2\} \cup Q) - M(\Cscr \cup \{S\} \cup Q)$
    \If{$\Delta M > 0$}
        \State $Q \gets Q \cup \{S_1, S_2\}$
    \Else
        \State $\Cscr \gets \Cscr \cup \{S\}$
    \EndIf
\EndWhile
\State $\Cscr \gets \Cscr \cup Q$
\State \Return $\Cscr$
\end{algorithmic}
\end{algorithm}

\paragraph{Bipartition step.}
The procedure $\textsc{BipartitionQIGNN}(G[S])$ applies the
bipartition QIGNN of Section~\ref{sec:method_qubo_k2} to the
induced subgraph $G[S]$. Several independent runs are performed
with different random initializations and the run achieving the
best discrete subgraph modularity is kept. The result is a
partition of $S$ into two non-empty subsets $S_1$ and $S_2$.

\paragraph{Stopping criterion.}
A candidate split is accepted only when it strictly increases the
modularity of the global partition. This is a direct adaptation of
Newman's stopping rule for hierarchical modularity
maximization~\cite{newman2006modularity}: a split that does not
improve modularity at the global level is not retained, even if it
locally separates the subgraph. The criterion automatically
determines a suitable number of communities; the parameter
$k_{\max}$ acts only as a safety bound and can be set to infinity
when no upper limit is desired.

\paragraph{Subgraph induction.}
Each call to $\textsc{BipartitionQIGNN}$ operates on the induced
subgraph $G[S]$, which preserves the internal connectivity of $S$
but discards edges to the rest of the graph. The modularity matrix
$\B$ is recomputed for $G[S]$ at each level, with degrees and the
total edge count taken from the subgraph. The hierarchical
procedure therefore optimizes the modularity of $G$ globally while
performing only local QIGNN training at each step, which keeps
the per-step cost low even when the graph is large.

\paragraph{Properties.}
The algorithm has three notable properties. First, in the
balanced-split case the depth of the recursion is
$\Theta(\log_2 k)$ where $k$ is the final number of communities,
while the worst case (pathologically unbalanced splits) is
$O(|V|)$. In practice the modularity-gain criterion terminates
recursion early on graphs whose true number of communities is
moderate, and the observed maximum depth in the experiments of
Section~\ref{sec:exp_hierarchical} never exceeds five (attained
on Email-EU-core). Second, the algorithm produces an interpretable
hierarchy of communities as a by-product, which is useful for
exploratory analysis even when only the final flat partition is
needed. Third, because every internal step uses the stable
bipartition formulation~\eqref{eq:qubo_k2}, the overall
procedure is empirically deterministic: in the experiments of
Section~\ref{sec:exp_hierarchical}, the standard deviation of
modularity across five random seeds is below $10^{-4}$ on ten of
thirteen benchmarks~--~at the numerical noise floor~--~and at
most $5\cdot10^{-3}$ on the remaining three.

\newpage
